\chapter{Background and Related Work}
This chapter reviews the main approaches to the control of legged and
wheeled-legged robots, with particular attention to the combination of
model-based control and reinforcement learning. The discussion begins with the
characteristics of these robotic platforms and the challenges introduced by
dynamic locomotion, intermittent contacts, and underactuated floating-base
dynamics.\\
\
Model-based methods are then examined through reduced-order models, Model
Predictive Control, and Whole-Body Control. These techniques provide structured
and interpretable control solutions, but their performance depends on model
accuracy and is constrained by the computational cost of online optimization.
Learning-based locomotion offers a complementary approach, allowing policies to
adapt to complex dynamics and uncertain interactions through experience.\\
\
The chapter subsequently considers hybrid architectures that combine MPC and
Reinforcement Learning. Particular attention is given to hierarchical,
residual, and adaptive formulations, which assign different roles to the
learned component and the nominal controller. Finally, recent developments in
GPU-accelerated optimal control are discussed, since they make it possible to
execute many MPC instances in parallel and train reinforcement-learning
policies with the optimizer directly in the control loop.

\section{Related Works}
\label{sec:related_works}

Optimization-based control has become a standard approach to dynamic legged
locomotion. It formulates motion generation as an Optimal Control Problem in
which the robot dynamics, contact forces, physical constraints, and task
objectives are considered explicitly. Solving the complete problem online is
difficult because legged robots have many degrees of freedom and interact with
the environment through intermittent and unilateral contacts
\cite{legged_optimization}. Existing controllers therefore differ in the
dynamic model they employ, the way they represent contacts, and the division of
tasks between predictive planning and low-level control.\\
\
Reduced-order models make online optimization tractable by retaining the
variables that dominate balance and locomotion while neglecting part of the
articulated dynamics. The convex MPC developed for the MIT Cheetah 3 represents
the trunk as a single rigid body and optimizes the ground reaction forces over
a finite horizon \cite{cheetah_mpc}. Whole-body nonlinear MPC instead includes
the articulated dynamics and optimizes robot motion and contact interactions
within a single problem \cite{whole_body_nmpc}. This richer representation
reduces the separation between planning and control but results in a more
demanding optimization. Other work has addressed the numerical treatment of
three-dimensional orientation \cite{representation_free_mpc} and adapted solver
accuracy to the time constraints of whole-body locomotion control
\cite{fast_whole_body_mpc}.\\
\
Wheeled-legged robots introduce rolling contacts and additional actuated
degrees of freedom. Early experiments on ANYmal integrated wheel motion and
non-holonomic rolling constraints into a whole-body planning and control
framework \cite{keep_rollin}. Online trajectory optimization was later used to
produce hybrid walking-driving motions on rough terrain
\cite{rolling_deep}. A subsequent whole-body MPC formulation jointly optimized
torso motion, wheel motion, and contact forces while adapting the contact
sequence online \cite{wheeled_whole_body_mpc}. These works showed that the
wheels cannot be treated as an independent drivetrain. Their motion must be
coordinated with the leg configuration, body dynamics, and contact forces.\\
\
Research on wheeled-legged systems has mainly addressed terrain traversal and
the transition between rolling and stepping. At higher speeds, vehicle-like
effects become increasingly relevant. Acceleration and cornering redistribute
the normal loads among the wheels, reducing the available contact margin and
increasing the risk of slip or rollover. A recent controller for the Unitree
Go2-W combined an MPC planner for active roll regulation with a learned
low-level policy. The legs were used as an active suspension system to lean the
body into turns and reduce lateral load transfer during physical racing
experiments \cite{racing}.\\
\
Reinforcement Learning provides a different approach to the uncertainties and
contact interactions involved in locomotion. Policies trained in simulation
have been transferred to physical quadrupeds for agile movements and locomotion
over challenging terrain \cite{agile_motor_skills,challenging_terrain}.
Perceptive policies have also demonstrated locomotion in natural environments
using proprioceptive and exteroceptive information
\cite{perceptive_rl}. These methods can learn behaviours that are difficult to
describe through an analytical model, but their performance depends on the
training distribution, reward formulation, observation design, and simulation
accuracy.\\
\
The computational structure of RL has encouraged massively parallel
simulation. Thousands of robots can be simulated simultaneously on a GPU,
producing large batches of experience for policy optimization. This approach
has substantially reduced the time required to train locomotion policies and
supports curriculum learning across many parallel environments
\cite{parallel_rl}. A purely model-free policy, however, does not explicitly
encode the physical constraints or the structure available to an
optimization-based controller. This limitation has motivated methods that
combine learning and model-based control.\\
\
One family of hybrid methods uses model-based optimization to generate
references or guidance for a learned policy. Deep Tracking Control computes
reference motions through trajectory optimization and trains an RL policy to
track them, while retaining the planner during deployment
\cite{deep_tracking}. A related approach uses a reduced-order model to predict
desired footholds and lets RL realize those footholds without tracking a
complete reference trajectory \cite{footstep_rl}. In both cases, the model
guides the learned behaviour, while the policy remains responsible for
producing the low-level control action.\\
\
Hierarchical methods assign planning and execution to different components.
The policy may select footholds, contact timings, gait parameters, or
navigation commands, while MPC computes a dynamically consistent motion. In
\cite{non_gaited}, a high-level RL policy selects acyclic contact patterns and
navigation commands for a low-level MPC. This removes the combinatorial search
over contact sequences from the continuous optimization and allows the same
architecture to address legged and hybrid wheeled-legged locomotion.\\
\
A separate family uses learning to modify the model or parameters employed by
the optimizer. Learned corrections have been introduced to reduce the
discrepancy between reduced-order predictions and the full-order dynamics of a
quadruped \cite{model_gap}. Modular residual methods apply learned corrections
to individual parts of a model-based controller, including its dynamics model
and footstep planner \cite{modular_residual}. In the adaptive architecture
presented in \cite{adaptive_mpc}, RL modifies the predictive dynamics,
swing-foot parameters, and gait frequency used by MPC. The learned action
therefore changes the optimization problem rather than directly correcting the
actuator command.\\
\
Residual policy learning adopts a more direct decomposition. A conventional
controller supplies a nominal action, while RL learns only the correction
needed to improve the resulting behaviour. This principle was introduced for
contact-rich manipulation tasks, where a feedback controller captured the
known structure and a residual policy compensated for effects that were
difficult to model \cite{residual_robot_control}. In locomotion, a residual can
be applied to the dynamic model, a planned reference, a joint setpoint, or a
torque command. These alternatives differ in the authority assigned to the
policy and in how much of the nominal controller remains responsible for
physical consistency.\\
\
The work most closely related to this thesis is Residual MPC, which executes an
MPC controller and a residual policy concurrently and combines their outputs at
the joint-control level \cite{residual_mpc}. The MPC provides a structured
control prior, while the policy learns corrections for modelling errors and
behaviours that are difficult to encode in the optimization. The paper compares
different residual action spaces and blending strategies, showing that the
interface between the two controllers affects both learning and final
locomotion behaviour. Unlike approaches based on offline demonstrations, MPC
remains active during training and deployment.\\
\
This integration creates a substantial computational requirement. An MPC
problem must be solved for every simulated robot at each control update.
Optimal-control solvers for legged locomotion have traditionally been designed
for sequential execution on CPUs, whereas modern RL environments simulate
thousands of agents in parallel on GPUs. The resulting separation between CPU
optimization and GPU training limits the number of MPC instances that can be
kept inside the learning loop \cite{residual_mpc,mpx}.\\
\
CusADi addressed this problem by translating symbolic expressions into
GPU-compatible operations and evaluating large batches of approximate
optimal-control problems in parallel \cite{cusadi}. ReLU-QP followed a
different direction by reformulating an ADMM-based quadratic-programming solver
using operations efficiently supported by machine-learning frameworks and GPUs
\cite{relu_qp}. MPX extends GPU acceleration to a primal-dual iLQR formulation
implemented in JAX. Associative scans parallelize computations along the
prediction horizon, while batching supports many independent MPC instances on
the same device \cite{mpx}. This enables MPC to remain directly inside a
massively parallel RL training loop.\\
\
The literature therefore combines model-based control and learning at several
different levels. The model may generate demonstrations, supply reference
motions, determine contact sequences, define the dynamics used for prediction,
or produce the nominal actuator command. In turn, the learned policy may track
a model-based plan, select high-level decisions, adapt the optimization model,
or provide a residual correction. This thesis focuses on the last relation:
the model-based controller supplies the nominal behaviour, while RL learns a
correction through interaction with the simulated system. The exact interface
between the two components depends on the robot platform and is described in
the architecture chapter.

\section{Legged and Wheeled-Legged Locomotion}
\label{sec:legged_wheeled_legged_locomotion}

Legged robots move through the environment directly through force contacts
exchange with the surrounding terrain. Unlike conventional wheeled
vehicles, they do not require a continuous and approximately flat path. Each
foot can be placed on a separate support surface, allowing the robot to cross
steps, gaps, debris, and other irregular terrain. This mobility makes legged
systems suitable for environments in which fixed infrastructure is unavailable
or ordinary vehicles cannot operate \cite{legged_optimization}.\\
\
This capability introduces a demanding control problem. A legged robot is a
floating-base mechanical system whose body is not directly actuated. Its motion
must be generated indirectly through the forces exchanged at the active
contacts. These contacts change during locomotion and are governed by
unilateral and friction constraints: a foot can push against the terrain but
cannot pull on it, while excessive tangential force may cause slipping. The
controller must coordinate joint motion, contact timing, and ground reaction
forces while maintaining balance and tracking the desired movement
\cite{legged_optimization}.\\
\
The complexity of locomotion also depends on the selected gait. Periodic gaits,
such as walking, trotting, and bounding, define recurring sequences of stance
and swing phases. During stance, the contacting limbs support and accelerate
the body. During swing, the remaining limbs are repositioned for the next
contact. Faster motions reduce the time available for feedback and may include
phases with a small support region or no ground contact. Aperiodic locomotion
provides more freedom to react to the terrain, but selecting contact timings
while planning the continuous motion makes the control problem considerably
harder \cite{legged_optimization,non_gaited}.\\
\
Model uncertainty further separates the predicted and actual motion. Impacts,
actuator limitations, delayed measurements, uncertain friction, and imperfect
terrain information cannot always be represented accurately. These effects are
particularly important when a controller relies on a reduced-order model, since
the neglected articulated dynamics still influence the physical robot. Robust
locomotion therefore requires continuous feedback and repeated adaptation of
the planned motion.\\
\
Wheeled-legged robots extend this mobility by placing wheels at the ends of
some or all limbs. On regular terrain, rolling avoids the repeated
acceleration and deceleration associated with conventional stepping and can
provide faster and more energy-efficient locomotion
\cite{wheeled_legged_survey,keep_rollin}. When the terrain becomes
discontinuous or unsuitable for rolling, the articulated legs can adjust the
body posture, negotiate obstacles, or initiate a stepping motion
\cite{rolling_deep}. The same platform can therefore use rolling for efficient
traversal and leg motion when the geometry of the environment requires it.\\
\
The mechanical arrangement of the wheels determines the available mobility.
Some platforms use steerable wheels, whereas others mount non-steerable wheels
at the ends of the legs. In the latter case, each contact can move along the
rolling direction but should not slide laterally. The controller must enforce
this non-holonomic condition while coordinating wheel velocity, leg
configuration, body motion, and contact forces
\cite{keep_rollin,wheeled_whole_body_mpc}. The wheels increase the set of
feasible motions, but also couple locomotion and posture control more tightly
than in a purely legged system.\\
\
Wheeled-legged locomotion can be divided into rolling, stepping, and hybrid
motions. During rolling, the wheels remain in contact with the ground and the
legs regulate the posture of the body. During stepping, one or more wheel
contacts are lifted and repositioned as ordinary feet. Hybrid locomotion
combines these modes within the same manoeuvre, allowing the robot to drive
across traversable regions and step over local obstacles
\cite{rolling_deep}. The contact sequence may be prescribed in advance or
adapted online according to the state of the robot and its kinematic limits
\cite{wheeled_whole_body_mpc}.\\
\
The addition of wheels does not remove the difficulties associated with legged
locomotion. Rolling contacts introduce further constraints, and transitions
between driving and stepping modify the support configuration. At higher
speeds, acceleration and cornering also redistribute the normal loads between
the contacts. This lateral load transfer reduces the available contact margin
and can increase the risk of wheel slip or rollover. On skid-steered platforms,
yaw is generated through differential wheel torques, further coupling
longitudinal motion, lateral dynamics, and body posture \cite{racing}.\\
\
Articulated legs can be used to mitigate these effects. By changing the
relative extension of the legs on either side of the body, a wheeled-legged
robot can regulate its roll angle and redistribute the contact loads. Physical
experiments on the Unitree Go2-W demonstrated this principle by using the legs
as an active suspension system and leaning the body into turns
\cite{racing}. This extends the role of the legs beyond obstacle traversal:
their configuration also affects stability and vehicle-like dynamics during
fast rolling.\\
\
Legged and wheeled-legged robots thus share the need to regulate an
underactuated floating body through constrained ground contacts. Wheels improve
continuous mobility on suitable terrain, but introduce rolling constraints,
additional actuation, and high-speed effects that are absent or less prominent
in conventional stepping. These properties motivate controllers that can plan
body motion and contact forces online while remaining sufficiently efficient
for closed-loop operation.

\section{Model-Based Locomotion Control}

\subsubsection{Double Forces Constrained Inverted Pendulum MPC}
\label{subsubsec:dfcip_mpc}

The predictive controller used for TITA is based on the Double Forces
Constrained Inverted Pendulum (DFCIP). The model describes the centre of mass
as a point mass supported by two ground reaction forces, one for each wheel.
The motion of the wheeled support is represented by a virtual unicycle located
at the midpoint between the wheel contacts.\\
\
The DFCIP state is
%
\begin{equation}
    \mathbf{x}
    =
    \begin{bmatrix}
        \mathbf{p}_{\mathrm{CoM}}^{\top} &
        \mathbf{v}_{\mathrm{CoM}}^{\top} &
        \mathbf{c}^{\top} &
        v_{c,z} &
        \theta &
        v &
        \omega
    \end{bmatrix}^{\top}
    \in\mathbb{R}^{13},
    \label{eq:dfcip_state}
\end{equation}
%
where $\mathbf{p}_{\mathrm{CoM}}$ and $\mathbf{v}_{\mathrm{CoM}}$ are the
centre-of-mass position and velocity. The vector $\mathbf{c}$ denotes the
midpoint between the two wheel contacts, while $v_{c,z}$ is its vertical
velocity. The heading angle, forward speed, and yaw rate of the virtual
unicycle are denoted by $\theta$, $v$, and $\omega$.\\
\
The control input is
%
\begin{equation}
    \mathbf{u}
    =
    \begin{bmatrix}
        a &
        a_{c,z} &
        \alpha &
        \mathbf{F}_{L}^{\top} &
        \mathbf{F}_{R}^{\top}
    \end{bmatrix}^{\top}
    \in\mathbb{R}^{9}.
    \label{eq:dfcip_input}
\end{equation}
%
These are the forward, angular and vertical acceleration
of the CoM, and the ground reaction forces applied to the wheel contact points.
\
The centre-of-mass dynamics are
%
\begin{equation}
    \dot{\mathbf{p}}_{\mathrm{CoM}}
    =
    \mathbf{v}_{\mathrm{CoM}},
    \qquad
    \dot{\mathbf{v}}_{\mathrm{CoM}}
    =
    \mathbf{g}
    +
    \frac{1}{m}
    \left(
        \mathbf{F}_{L}
        +
        \mathbf{F}_{R}
    \right).
    \label{eq:dfcip_com_dynamics}
\end{equation}
%
The midpoint between the wheel contacts follows the unicycle dynamics
%
\begin{equation}
\begin{aligned}
    \dot{c}_x
        &=
        v\cos\theta, &
    \dot{c}_y
        &=
        v\sin\theta, \\
    \dot{c}_z
        &=
        v_{c,z}, &
    \dot{v}_{c,z}
        &=
        0, \\
    \dot{\theta}
        &=
        \omega, &
    \dot{v}
        &=
        a, \qquad
    \dot{\omega}
        =
        \alpha.
\end{aligned}
\label{eq:dfcip_unicycle_dynamics}
\end{equation}
%
The equations for $\dot{c}_x$ and $\dot{c}_y$ enforce the nonholonomic
structure of the wheeled base: the support midpoint moves along the direction
defined by $\theta$ and has no independent lateral-velocity input. The model is
discretized using forward Euler integration with a prediction step of
$0.01\,\mathrm{s}$.\\
\
Let $d$ denote the distance between the wheels. Their contact positions are
reconstructed from the midpoint and heading:
%
\begin{equation}
\begin{aligned}
    \mathbf{p}_{c,L}
        &=
        \mathbf{c}
        +
        \mathbf{R}_{z}(\theta)
        \begin{bmatrix}
            0 & d/2 & 0
        \end{bmatrix}^{\top}, \\
    \mathbf{p}_{c,R}
        &=
        \mathbf{c}
        -
        \mathbf{R}_{z}(\theta)
        \begin{bmatrix}
            0 & d/2 & 0
        \end{bmatrix}^{\top}.
\end{aligned}
\label{eq:dfcip_contact_positions}
\end{equation}
%
Since the DFCIP does not include an explicit rotational state for the complete
robot, the consistency of the force distribution is expressed through the
moment-balance residual
%
\begin{equation}
    \mathbf{r}_{M,k}
    =
    \left(
        \mathbf{p}_{c,L,k}
        -
        \mathbf{p}_{\mathrm{CoM},k}
    \right)
    \times
    \mathbf{F}_{L,k}
    +
    \left(
        \mathbf{p}_{c,R,k}
        -
        \mathbf{p}_{\mathrm{CoM},k}
    \right)
    \times
    \mathbf{F}_{R,k}.
    \label{eq:dfcip_moment_residual}
\end{equation}
%
Penalizing this residual couples the location of the centre of mass to the
left--right force distribution. During turning, for example, the optimizer can
vary the tangential and vertical components of the two forces while limiting
the unmodelled resultant moment.\\
\
The DFCIP objective contains a larger set of possible tracking residuals. The stage cost is
%
\begin{equation}
\begin{aligned}
    \ell_k = 
        &w_{p_{xy}}
        \left(
            p_{\mathrm{CoM},xy,k}
            -p_{\mathrm{CoM},xy,k}^{\mathrm{ref}}
        \right)^2
        +w_{p_z}
        \left(
            p_{\mathrm{CoM},z,k}
            -p_{\mathrm{CoM},z,k}^{\mathrm{ref}}
        \right)^2
        \\
        &+w_{v_{xy}}
        \left(
            v_{\mathrm{CoM},xy,k}
            -v_{\mathrm{CoM},xy,k}^{\mathrm{ref}}
        \right)^2
        +w_{v_z}
        \left(
            v_{\mathrm{CoM},z,k}
            -v_{\mathrm{CoM},z,k}^{\mathrm{ref}}
        \right)^2
        \\
        &+w_c\left(c_k-c_k^{\mathrm{ref}}\right)^2
        +w_{v_c}
        \left(v_{c,z,k}-v_{c,z,k}^{\mathrm{ref}}\right)^2
        +w_{\theta}
        \left(\theta_k-\theta_k^{\mathrm{ref}}\right)^2
        \\
        &+w_v\left(v_k-v_k^{\mathrm{ref}}\right)^2
        +w_{\omega}\left(\omega_k-\omega_k^{\mathrm{ref}}\right)^2
        \\
        &+w_a\left(a_k-a_k^{\mathrm{ref}}\right)^2
        +w_{a_z}
        \left(a_{c,z,k}-a_{c,z,k}^{\mathrm{ref}}\right)^2
        +w_{\alpha}
        \left(\alpha_k-\alpha_k^{\mathrm{ref}}\right)^2
        \\
        &+\sum_{i\in\{L,R\}}
        \left[
            w_{F_{xy}}
            \left(F_{i,xy,k}-F_{i,xy,k}^{\mathrm{ref}}\right)^2
            +w_{F_z}
            \left(F_{i,z,k}-F_{i,z,k}^{\mathrm{ref}}\right)^2
        \right]
        \\
        &+w_{\mathrm{eq}}
        \left[
            v_{c,z,k}^{2}
            +r_{M,k}
            +\sum_{i\in\{L,R\}}
             \min\left(F_{i,z,k},0\right)^2
        \right]
\end{aligned}
\label{eq:dfcip_stage_cost}
\end{equation}
%
Where momentum is expressed as 
%
\begin{equation}
    r_{F_z,k}
    =
    \min\left(F_{L,z,k},0\right)^2
    +
    \min\left(F_{R,z,k},0\right)^2.
    \label{eq:dfcip_unilateral_force}
\end{equation}
%
It is zero when both vertical forces are nonnegative and increases when either
wheel force would pull on the ground. The residual $v_{c,z,k}^{2}$ keeps the
virtual wheel midpoint on the flat contact surface. These relations are
treated as soft constraints through the large coefficient
$w_{\mathrm{eq}}$.\\
\
The reference generator assigns
%
\begin{equation}
\begin{aligned}
    p_{\mathrm{CoM},z}^{\mathrm{ref}}
        &= 0.4\,\mathrm{m}, \\
    v_{\mathrm{CoM},z}^{\mathrm{ref}}
        &= 0, \\
    v_{c,z}^{\mathrm{ref}}
        &= 0, \\
    \mathbf{F}_{L}^{\mathrm{ref}}
        &=
        \begin{bmatrix}
            0 & 0 & mg/2
        \end{bmatrix}^{\top}, \\
    \mathbf{F}_{R}^{\mathrm{ref}}
        &=
        \begin{bmatrix}
            0 & 0 & mg/2
        \end{bmatrix}^{\top}.
\end{aligned}
\label{eq:dfcip_nominal_references}
\end{equation}
%
The desired forward speed and yaw rate are approached through
acceleration-limited profiles. At each prediction step,
%
\begin{equation}
\begin{aligned}
    v_{k+1}^{\mathrm{ref}}
        &=
        v_k^{\mathrm{ref}}
        +
        \operatorname{clip}
        \left(
            v^{\mathrm{cmd}}-v_k^{\mathrm{ref}},
            -a_{\max}\Delta t,
            a_{\max}\Delta t
        \right), \\
    \omega_{k+1}^{\mathrm{ref}}
        &=
        \omega_k^{\mathrm{ref}}
        +
        \operatorname{clip}
        \left(
            \omega^{\mathrm{cmd}}-\omega_k^{\mathrm{ref}},
            -\alpha_{\max}\Delta t,
            \alpha_{\max}\Delta t
        \right).
\end{aligned}
\label{eq:dfcip_velocity_ramp}
\end{equation}
%
The planar centre-of-mass velocity reference is aligned with the predicted
heading:
%
\begin{equation}
    \mathbf{v}_{\mathrm{CoM},xy,k}^{\mathrm{ref}}
    =
    v_k^{\mathrm{ref}}
    \begin{bmatrix}
        \cos\theta_k^{\mathrm{ref}} \\
        \sin\theta_k^{\mathrm{ref}}
    \end{bmatrix}.
    \label{eq:dfcip_com_velocity_reference}
\end{equation}
%
The objective therefore tracks forward speed in two related forms: through the
scalar unicycle speed $v$ and through the planar centre-of-mass velocity. The
first represents the motion of the wheel midpoint along its heading, whereas
the second regulates the global translational motion of the centre of mass.\\
\
At the terminal node, input-dependent terms are removed and an additional
stability residual is introduced:
%
\begin{equation}
    \mathbf{r}_{\mathrm{stab},N}
    =
    \mathbf{p}_{\mathrm{CoM},xy,N}
    -
    \mathbf{c}_{xy,N}.
    \label{eq:dfcip_terminal_stability}
\end{equation}
%
The terminal cost is
%
\begin{equation}
\begin{aligned}
    \ell_N =
        &w_{p_{xy}}
        \left(
            p_{\mathrm{CoM},xy,N}
            -p_{\mathrm{CoM},xy,N}^{\mathrm{ref}}
        \right)^2
        +w_{p_z}
        \left(
            p_{\mathrm{CoM},z,N}
            -p_{\mathrm{CoM},z,N}^{\mathrm{ref}}
        \right)^2
        \\
        &+w_{v_{xy}}
        \left(
            v_{\mathrm{CoM},xy,N}
            -v_{\mathrm{CoM},xy,N}^{\mathrm{ref}}
        \right)^2
        +w_{v_z}
        \left(
            v_{\mathrm{CoM},z,N}
            -v_{\mathrm{CoM},z,N}^{\mathrm{ref}}
        \right)^2
        \\
        &+w_c
        \left(
            c_N-c_N^{\mathrm{ref}}
        \right)^2
        +w_{v_c}
        \left(
            v_{c,z,N}-v_{c,z,N}^{\mathrm{ref}}
        \right)^2
        +w_{\theta}
        \left(
            \theta_N-\theta_N^{\mathrm{ref}}
        \right)^2
        \\
        &+w_v
        \left(
            v_N-v_N^{\mathrm{ref}}
        \right)^2
        +w_{\omega}
        \left(
            \omega_N-\omega_N^{\mathrm{ref}}
        \right)^2
        \\
        &+w_{\mathrm{eq}}
        \left[
            v_{c,z,N}^{2}
            +\left(
                p_{\mathrm{CoM},xy,N}-c_{xy,N}
            \right)^2
        \right]
\end{aligned}
\label{eq:dfcip_terminal_cost}
\end{equation}
%
The terminal condition encourages the centre-of-mass projection to return over
the midpoint of the two wheels. It acts as a reduced stability objective at the
end of the finite horizon. Moment balance and unilateral-force penalties are
stage terms and are not included in the implemented terminal cost.\\
\
The tasks solved directly by the TITA DFCIP-MPC are:

\begin{itemize}
    \item regulating the centre-of-mass $p_{\mathrm{CoM}}$;
    \item tracking the centre-of-mass velocity $v_{\mathrm{CoM}}$;
    \item tracking the forward speed of the wheeled support $v$;
    \item tracking the commanded yaw rate $\omega$;
    \item regularizing forward, vertical, and yaw accelerations $a$, $a_{c,z}$, $\alpha$;
    \item distributing the ground reaction forces between the two wheels $\mathbf{F}_L$, $\mathbf{F}_R$;
    \item regularizing tangential and vertical contact forces around their nominal references $F_{i,xy}$, $F_{i,z}$;
    \item enforcing moment balance through a soft residual $r_M$;
    \item maintaining zero vertical velocity of the wheel midpoint $v_{c,z}$;
    \item penalizing negative wheel normal forces $F_{i,z}<0$;
\end{itemize}
\
Planar centre-of-mass position, the position of the virtual contact midpoint,
its vertical-velocity reference, and the absolute heading angle appear in the
general objective implementation but have zero weight in the configuration
used for this work. They must therefore not be listed as active MPC tasks.\\
\
The current controller assumes permanent contact with flat ground. It does not
switch between stance and flight dynamics, and the contact forces are not set
to zero during any part of the prediction.\\
\
The DFCIP-MPC does not take into account robot dynamics and ground-reaction forces constraints. It only penalizes
negative vertical forces through
Equation~\eqref{eq:dfcip_unilateral_force}, these are handled by
the subsequent Whole-Body Controller \ref{subsubsec:wbc_dfcip}, which maps
the first MPC solution into torques while enforcing the constraints omitted from the reduced
model.\\
\
For TITA, the prediction horizon contains 50 intervals of
$0.01\,\mathrm{s}$, corresponding to $0.5\,\mathrm{s}$. The MPC is updated at
$100\,\mathrm{Hz}$. One solver iteration is performed at each update, following
a real-time iteration scheme, and the previous solution is shifted by one node
to initialize the next OCP. The same reduced-order problem is solved
independently for each simulated environment through the batched GPU
implementation provided by MPX \cite{mpx}.

\subsubsection{Single Rigid Body Dynamics MPC}
\label{subsubsec:srbd_mpc}

Single Rigid Body Dynamics (SRBD) provides a reduced representation of a legged
robot in which the complete multibody system is approximated as a single rigid
body, not including the limbs in the optimized state. Their
effect on the trunk is represented through the ground reaction forces applied
at the contact points. This approximation reduces the dimension and
nonlinearity of the Optimal Control Problem, making it suitable for
online locomotion control \cite{mpx}.\\
\
Let $\mathbf{p}\in\mathbb{R}^{3}$ denote the position of the centre of mass,
$\mathbf{q}\in\mathbb{S}^{3}$ the base orientation, and
$\mathbf{v},\boldsymbol{\omega}\in\mathbb{R}^{3}$ the linear and angular
velocities. The SRBD state is
%
\begin{equation}
    \mathbf{x}
    =
    \begin{bmatrix}
        \mathbf{p}^{\top} &
        \mathbf{q}^{\top} &
        \mathbf{v}^{\top} &
        \boldsymbol{\omega}^{\top}
    \end{bmatrix}^{\top}
    \in\mathbb{R}^{13}.
    \label{eq:srbd_state}
\end{equation}
%
For a quadruped with four point contacts, the control input contains the ground
reaction force associated with each foot:
%
\begin{equation}
    \mathbf{u}
    =
    \begin{bmatrix}
        \mathbf{F}_{FL}^{\top} &
        \mathbf{F}_{FR}^{\top} &
        \mathbf{F}_{HL}^{\top} &
        \mathbf{F}_{HR}^{\top}
    \end{bmatrix}^{\top}
    \in\mathbb{R}^{12}.
    \label{eq:srbd_input}
\end{equation}
%
The subscripts $FL$, $FR$, $HL$, and $HR$ identify the front-left,
front-right, hind-left, and hind-right feet. The control variables are therefore
contact forces rather than actuator torques.\\
\
A predefined contact schedule determines which forces act on the body at each
prediction node. Let $\gamma_{i,k}\in\{0,1\}$ indicate whether foot $i$ is in
stance at node $k$. The translational dynamics are
%
\begin{equation}
    m\dot{\mathbf{v}}_k
    =
    m\mathbf{g}
    +
    \sum_{i=1}^{n_c}
    \gamma_{i,k}\mathbf{F}_{i,k},
    \label{eq:srbd_linear_dynamics}
\end{equation}
%
where $m$ is the total mass and $n_c=4$ for Lite3. Forces associated with swing
feet are removed from the predicted body dynamics by the contact mask
$\gamma_{i,k}$.\\
\
The contact forces also generate moments about the centre of mass. Assuming a
constant centroidal inertia $\mathbf{I}$ expressed in the body frame, the
rotational dynamics are
%
\begin{equation}
    \mathbf{I}\dot{\boldsymbol{\omega}}_k
    +
    \boldsymbol{\omega}_k
    \times
    \mathbf{I}\boldsymbol{\omega}_k
    =
    \mathbf{R}(\mathbf{q}_k)^{\top}
    \sum_{i=1}^{n_c}
    \gamma_{i,k}
    \left(
        \mathbf{p}_{f,i,k}-\mathbf{p}_k
    \right)
    \times
    \mathbf{F}_{i,k},
    \label{eq:srbd_rotational_dynamics}
\end{equation}
%
where $\mathbf{p}_{f,i,k}$ is the position of the $i$-th foot and
$\mathbf{R}(\mathbf{q}_k)$ is the rotation from the body frame to the world
frame. The moment generated by each force depends on its lever arm with respect
to the centre of mass. The MPC can therefore regulate the orientation and
angular velocity of the body by redistributing the forces among the active
contacts. This force-distribution mechanism is also used in MPC formulations
for high-speed wheeled-legged locomotion \cite{racing}.\\
\
The discrete dynamics used in the implementation are obtained through
semi-implicit Euler integration. Linear and angular velocities are updated
first:
%
\begin{equation}
\begin{aligned}
    \mathbf{v}_{k+1}
        &=
        \mathbf{v}_k
        +
        \Delta t
        \left(
            \mathbf{g}
            +
            \frac{1}{m}
            \sum_{i=1}^{n_c}
            \gamma_{i,k}\mathbf{F}_{i,k}
        \right), \\
    \boldsymbol{\omega}_{k+1}
        &=
        \boldsymbol{\omega}_k
        +
        \Delta t\,
        \mathbf{I}^{-1}
        \left[
            \mathbf{R}(\mathbf{q}_k)^{\top}
            \sum_{i=1}^{n_c}
            \gamma_{i,k}
            \left(
                \mathbf{p}_{f,i,k}-\mathbf{p}_k
            \right)
            \times
            \mathbf{F}_{i,k}
            -
            \boldsymbol{\omega}_k
            \times
            \mathbf{I}\boldsymbol{\omega}_k
        \right].
\end{aligned}
\label{eq:srbd_velocity_update}
\end{equation}
%
The position and quaternion are then propagated using the updated velocities:
%
\begin{equation}
\begin{aligned}
    \mathbf{p}_{k+1}
        &=
        \mathbf{p}_k
        +
        \Delta t\,\mathbf{v}_{k+1}, \\
    \mathbf{q}_{k+1}
        &=
        \mathbf{q}_k
        \oplus
        \left(
            \Delta t\,\boldsymbol{\omega}_{k+1}
        \right),
\end{aligned}
\label{eq:srbd_pose_update}
\end{equation}
%
where $\oplus$ denotes quaternion integration.\\
\
The SRBD stage cost used for Lite3 is
%
\begin{equation}
\begin{aligned}
    \ell_k ={}&
    \left\|
        \mathbf{p}_k-\mathbf{p}^{\mathrm{ref}}_k
    \right\|_{\mathbf{Q}_p}^{2}
    +
    \left\|
        \mathbf{e}_{R}
        \left(
            \mathbf{q}_k,\mathbf{q}^{\mathrm{ref}}_k
        \right)
    \right\|_{\mathbf{Q}_R}^{2}
    \\
    &+
    \left\|
        \mathbf{v}_k-\mathbf{v}^{\mathrm{ref}}_k
    \right\|_{\mathbf{Q}_v}^{2}
    +
    \left\|
        \boldsymbol{\omega}_k
        -
        \boldsymbol{\omega}^{\mathrm{ref}}_k
    \right\|_{\mathbf{Q}_{\omega}}^{2}
    \\
    &+
    \sum_{i=1}^{n_c}
    \left\|
        \mathbf{F}_{i,k}
    \right\|_{\mathbf{Q}_F}^{2}
    +
    \sum_{i=1}^{n_c}
    \phi_{\mu}
    \left(
        \mathbf{F}_{i,k}
    \right),
\end{aligned}
\label{eq:lite3_srbd_stage_cost}
\end{equation}
%
where
$\|\mathbf{z}\|_{\mathbf{A}}^{2}
=\mathbf{z}^{\top}\mathbf{A}\mathbf{z}$ and
$\mathbf{e}_{R}$ is the three-dimensional orientation error computed from the
current and reference quaternions. The terminal cost contains the four state
tracking terms but omits force regularization and the friction penalty.\\
\
The MPC regulates the vertical position of the body, but does not track an absolute
planar position. Similarly, it penalizes roll and pitch errors but does not
track an absolute yaw angle. Planar movement and turning are instead specified
through linear- and angular-velocity references.\\
\
The friction term in Equation~\eqref{eq:lite3_srbd_stage_cost} is based on the
smoothed residual
%
\begin{equation}
    r_{\mu,i,k}
    =
    \mu F_{i,z,k}
    -
    \sqrt{
        F_{i,x,k}^{2}
        +
        F_{i,y,k}^{2}
        +
        \varepsilon
    },
    \label{eq:lite3_friction_residual}
\end{equation}
%
with $\mu=0.7$ and $\varepsilon=10^{-2}$ in the current implementation. A
relaxed barrier $\phi_{\mu}$ increases the cost when this residual approaches
or crosses zero. The friction cone is therefore treated as a soft constraint
inside the OCP.\\
\
The tasks solved directly by the Lite3 SRBD-MPC are:

\begin{itemize}
    \item tracking the commanded forward and lateral base velocities;
    \item regulating the vertical base velocity around zero;
    \item maintaining the desired base height;
    \item stabilizing the roll and pitch orientation of the trunk;
    \item regulating the roll and pitch angular velocities around zero;
    \item tracking the commanded yaw rate;
    \item distributing the ground reaction forces among the scheduled stance
          contacts;
    \item generating the resultant force required for linear body motion;
    \item generating the resultant moment required for rotational body motion;
    \item regularizing the magnitude of the optimized ground reaction forces;
    \item discouraging contact forces outside the friction cone.
\end{itemize}

Absolute planar position and absolute yaw orientation are not active tasks
because their weights are zero. Joint positions, joint velocities, actuator
torques, and swing-foot tracking are also absent from the OCP. They cannot be
optimized directly because they are not included in the SRBD state.\\
\
The velocity references and contact parameters are produced before solving the
OCP. The commanded planar velocity is rotated into the world frame using the
current heading:
%
\begin{equation}
    \mathbf{v}_{xy}^{\mathrm{ref}}
    =
    \mathbf{R}_{z}(\psi)
    \begin{bmatrix}
        v_x^{\mathrm{cmd}} \\
        v_y^{\mathrm{cmd}}
    \end{bmatrix}.
    \label{eq:lite3_world_velocity_reference}
\end{equation}
%
This velocity is integrated over the horizon to obtain the planar position
reference:
%
\begin{equation}
    \mathbf{p}_{xy,k}^{\mathrm{ref}}
    =
    \mathbf{p}_{xy,0}
    +
    k\Delta t\,
    \mathbf{v}_{xy}^{\mathrm{ref}}.
    \label{eq:lite3_position_reference}
\end{equation}
%
Although this reference is generated, it does not contribute directly to the
current objective because the planar entries of $\mathbf{Q}_p$ are zero. The
base height reference remains constant, and the reference orientation is
upright when terrain estimation is disabled. The angular-velocity reference is
%
\begin{equation}
    \boldsymbol{\omega}^{\mathrm{ref}}
    =
    \begin{bmatrix}
        0 &
        0 &
        \omega_z^{\mathrm{cmd}}
    \end{bmatrix}^{\top}.
    \label{eq:lite3_angular_velocity_reference}
\end{equation}
%
A separate gait generator specifies the contact schedule. Lite3 uses a trot in
which the front-left and hind-right legs share the same phase, while the
front-right and hind-left legs form the opposite pair. The foothold selected
for each swing leg follows a velocity-based rule of the form
%
\begin{equation}
    \mathbf{p}_{f,i,xy}^{\mathrm{target}}
    =
    \mathbf{p}_{f,i,xy}^{0}
    +
    \frac{D}{2f_{\mathrm{step}}}
    \mathbf{v}_{xy}^{\mathrm{ref}}
    +
    \sqrt{\frac{h}{g}}
    \left(
        \mathbf{v}_{xy}
        -
        \mathbf{v}_{xy}^{\mathrm{ref}}
    \right),
    \label{eq:lite3_foothold_selection}
\end{equation}
%
where $D$ is the duty factor, $f_{\mathrm{step}}$ is the step frequency, and
$h$ is the desired base height. Smooth polynomial trajectories connect the
lift-off location to the target foothold and raise the foot by the prescribed
step height. Early measured contacts are used to stop the swing trajectory.\\
\
Foothold selection and swing-trajectory generation must not be described as
tasks solved by the MPC. The generated contact locations enter the SRBD
dynamics as moment arms, whereas the swing-foot position and velocity
references are passed to the low-level controller. For a stance leg, this
controller maps the first optimized ground reaction force into joint torques
through the foot Jacobian. For a swing leg, it combines Cartesian feedback with
a dynamics-based feedforward term to track the generated foot trajectory.
Thus, contact scheduling determines where forces may act, the MPC computes
their required values, and the low-level controller realizes the resulting
stance and swing behavior.\\
\
For Lite3, the prediction horizon consists of 25 intervals of
$0.02\,\mathrm{s}$, corresponding to $0.5\,\mathrm{s}$. The MPC is updated at
$50\,\mathrm{Hz}$, while the low-level controller runs at $200\,\mathrm{Hz}$.
The previous state and control trajectories are shifted and reused as the
initial guess at the next MPC update.


\subsection{Whole-Body Control}
\label{subsubsec:wbc_dfcip}

\section{Learning-Based Locomotion Control}

\subsection{End-to-End Reinforcement Learning}

\subsection{Residual Reinforcement Learning}

\section{Hybrid MPC and Reinforcement Learning Architectures}

\subsection{Hierarchical Architectures}

\subsection{Residual Architectures}

\subsection{Adaptive MPC Architectures}

\section{GPU-Accelerated Optimization and Learning}

\section{Positioning of This Thesis}